\begin{trapbox}

	\begin{description}[style=nextline, leftmargin=2.8em, labelsep=0.5em, itemsep=0.35em]
		\item[\textbf{\textcolor{CTrap}{易错点一（忽略相交线）}}]
			认为直线垂直平面内两条直线（即使平行）即可得线面垂直。
		\item[\textbf{\textcolor{CTrap}{正确理解（两线必须相交）：}}]
			必须垂直两条相交直线才能推出线面垂直；若两线平行，则不能保证线面垂直。
		\item[\textbf{\textcolor{CTrap}{错因分析（记忆遗漏关键）：}}]
			对线面垂直判定定理的记忆不完整，忽略了“相交”这一关键条件。
	\end{description}

	\medskip
	\begin{description}[style=nextline, leftmargin=2.8em, labelsep=0.5em, itemsep=0.35em]
		\item[\textbf{\textcolor{CTrap}{易错点二（面面垂直性质误用）}}]
			已知两个平面垂直，直接认为一个平面内的任意直线都垂直于另一个平面。
		\item[\textbf{\textcolor{CTrap}{正确理解（需垂直于交线）：}}]
			必须是在一个平面内垂直于交线的直线才垂直于另一个平面；其他直线不一定垂直。
		\item[\textbf{\textcolor{CTrap}{错因分析（性质条件遗漏）：}}]
			面面垂直性质定理要求“垂直于交线”，学生常忽略这个条件。
	\end{description}

	\medskip
	\begin{description}[style=nextline, leftmargin=2.8em, labelsep=0.5em, itemsep=0.35em]
		\item[\textbf{\textcolor{CTrap}{易错点三（判定性质混淆）}}]
			在证明面面垂直时，试图直接通过两条直线垂直（线线垂直）来判定面面垂直，而非通过线面垂直。
		\item[\textbf{\textcolor{CTrap}{正确理解（判定需线面垂直）：}}]
			面面垂直的判定必须要求一个平面包含另一个平面的垂线，即先有线面垂直。
		\item[\textbf{\textcolor{CTrap}{错因分析（定理结构不清）：}}]
			混淆了线面垂直与面面垂直的判定条件，不知道面面垂直判定需要线面垂直作为中间步骤。
	\end{description}

\end{trapbox}
